\subsection*{Pipelined CG (Chronopoulos--Gear 1989): Implementation and Evaluation}

\paragraph{Algorithm.}
Standard unpreconditioned CG requires two sequential dot products per iteration:
one to compute $\alpha_k = \gamma_k / \delta_k$ (where $\gamma_k = \langle r_k, r_k \rangle$
and $\delta_k = \langle p_k, Ap_k \rangle$), and one to compute
$\beta_k = \gamma_{k+1} / \gamma_k$.
Each dot product is a global reduction that, in a distributed-memory (MPI) setting,
requires an \texttt{AllReduce} communication call.
Two sequential \texttt{AllReduce} calls per iteration are a significant bottleneck
when the communication latency is large.

The pipelined CG variant of Chronopoulos and Gear~\cite{chronopoulos1989s}
eliminates the sequential dependency by computing three dot products simultaneously:
\[
  \gamma_k = \langle r_k, r_k \rangle, \qquad
  \delta_k = \langle p_k, Ap_k \rangle, \qquad
  \eta_k   = \langle Ap_k, Ap_k \rangle,
\]
and deriving $\beta_k$ via the recurrence
\begin{equation}\label{eq:beta_recurrence}
  \beta_k = \frac{\alpha_k^2 \eta_k - \gamma_k}{\gamma_k},
\end{equation}
which predicts $\gamma_{k+1} / \gamma_k$ without requiring the residual $r_{k+1}$
to be computed first.
The derivation of~\eqref{eq:beta_recurrence} relies on the identity
$\langle r_k, Ap_k \rangle = \delta_k$,
which follows from $A$-conjugacy of consecutive search directions
($\langle p_{k-1}, Ap_k \rangle = 0$) and holds in exact arithmetic.

In the MPI setting, the three dot products can be communicated in a single
\texttt{AllReduce} call (three values), reducing the per-iteration communication
cost from two sequential \texttt{AllReduce} calls to one.

\paragraph{GPU implementation.}
We implemented \texttt{DevPCGSolver} in the \texttt{ngscuda} extension,
following the lecture-notes formulation (Algorithm~\ref{alg:pipcg}) exactly.
The solver uses a CUDA WHILE graph to capture the entire iteration loop on-device.

A key implementation challenge arose from the structure of the NGSolve assembled
matrix: the full stiffness matrix $A$ retains non-zero entries in the rows
corresponding to constrained (Dirichlet) degrees of freedom.
Applying $Ap = A p$ with $p$ in the free-dof space therefore produces a vector
$Ap$ with non-zero constrained components, which contaminate the residual
$r \leftarrow r - \alpha Ap$ over successive iterations and cause divergence.

The fix is to project $Ap$ back into the free-dof space after each matvec:
\[
  \tilde{s}_k = P(Ap_k), \quad \text{where } P = \text{Projector(FreeDofs)},
\]
and use $\tilde{s}_k$ in place of $Ap_k$ throughout.
This ensures $r_k$ remains in the free-dof space without requiring an explicit
preconditioner application to $r$ inside the loop.

\paragraph{Fused triple-dot kernel.}
Three separate cuBLAS dot-product calls each read two vectors, totalling
six vector reads.  We implemented a custom CUDA kernel that computes
$\gamma_k$, $\delta_k$, and $\eta_k$ in a single pass over the data,
reading $r$, $p$, and $Ap$ once each (three vector reads).
This halves the memory traffic for the reduction step.

\paragraph{Numerical stability.}
An early implementation that applied the formula~\eqref{eq:beta_recurrence}
to the full (unprojected) system diverged after many iterations.
This was initially misattributed to loss of $A$-conjugacy in floating-point arithmetic.
A control experiment using a NumPy reference implementation on the correctly
projected free-dof submatrix showed that~\eqref{eq:beta_recurrence} is
\emph{numerically stable} on this problem: the pipelined solver converged in
136 iterations (versus 135 for standard CG) with solution norms agreeing to
12 significant digits.
The divergence observed on GPU was caused entirely by the constrained-dof
contamination described above, not by the recurrence formula itself.

\paragraph{Performance results.}
Table~\ref{tab:pipcg} shows timing results for \texttt{DevCGSolver} (standard
preconditioned CG with WHILE graph) and \texttt{DevPCGSolver} (pipelined CG
with WHILE graph and fused triple-dot kernel) on four mesh sizes.

\begin{table}[h]
\centering
\begin{tabular}{rrrrrr}
\hline
$n$ & CG iters & CG (ms) & PipCG iters & PipCG (ms) & Speedup \\
\hline
 1\,295 & 102 & 3.13 & 100 & 3.09 & 1.01$\times$ \\
 2\,318 & 135 & 4.09 & 136 & 4.13 & 0.99$\times$ \\
 9\,918 & 204 & 6.27 & 212 & 7.43 & 0.85$\times$ \\
18\,249 & 280 & 9.01 & 284 & 9.28 & 0.97$\times$ \\
\hline
\end{tabular}
\caption{DevCGSolver vs DevPCGSolver on 3-D Poisson, H1 order 2, H100 GPU.}
\label{tab:pipcg}
\end{table}

\texttt{DevPCGSolver} achieves no measurable speedup over \texttt{DevCGSolver}.

\paragraph{Attribution of cost.}
Two candidate explanations were considered:
(a) the projector application on $Ap$ adds a kernel call that offsets the
fused-dot savings; and
(b) the per-iteration runtime is dominated by SpMV bandwidth, making
reduction savings invisible.

To separate these, a control experiment was run on a reaction-diffusion
problem with no Dirichlet boundary conditions, where all degrees of freedom
are free and the projector is the identity.  The relative performance was
identical (0.98$\times$--1.00$\times$), ruling out the projector as the
dominant cost.

The correct attribution is therefore \textbf{bandwidth dominance}:
the sparse matrix-vector product and preconditioner applications saturate
the H100 HBM3 memory bandwidth, and the savings from fewer intra-iteration
graph edges and reduced dot-product memory traffic are below the noise floor
of these dominant costs.

\paragraph{Remark on intra-iteration synchronisation.}
Inside a CUDA WHILE graph, host--device synchronisations are eliminated
entirely.  However, intra-iteration synchronisation still exists as
\emph{graph edges} encoding data dependencies between nodes.
Pipelined CG reduces the number of required edges (fewer reduction nodes)
and therefore allows potentially better GPU scheduling.
The reason no speedup is observed is not that this benefit is absent in
principle, but that it is dwarfed by SpMV bandwidth costs at the problem
sizes tested.

\paragraph{Multi-GPU outlook.}
The pipelining benefit becomes meaningful in a multi-GPU (MPI + CUDA) setting.
Standard CG requires two sequential \texttt{AllReduce} calls per iteration;
pipelined CG requires one (communicating $\gamma_k$, $\delta_k$, $\eta_k$
simultaneously).
On an H100 node with NVLink, \texttt{AllReduce} latency is of order
$10$--$50\,\mu\mathrm{s}$.
For 135 iterations, halving the \texttt{AllReduce} count saves
approximately $1$--$7\,\mathrm{ms}$ --- significant relative to a total
solve time of $4$--$9\,\mathrm{ms}$ on a single GPU.
Extending \texttt{DevPCGSolver} to multi-GPU via NCCL-based collective
operations is a natural direction for future work.
